% ---- score2 standard preamble (from: python3 score2_tools.py preamble) ----
\documentclass[10pt]{extarticle}
\usepackage[a4paper,margin=1.5cm]{geometry}
\usepackage[T1]{fontenc}
\usepackage{amsmath,amssymb}
\usepackage{graphicx}
\usepackage[export]{adjustbox}
\usepackage{multicol}
\usepackage{enumitem}
\usepackage{booktabs,array,tabularx}
\setlength{\columnsep}{0.7cm}
\setlength{\parindent}{0pt}
\setlength{\parskip}{3pt}
\pagestyle{empty}
\makeatletter
% figure registry (figdefs.tex is generated by the extract step)
\newcommand{\qfigdef}[4]{\@namedef{qf@file@#1}{#2}\@namedef{qf@w@#1}{#3}\@namedef{qf@wide@#1}{#4}}
\newcommand{\qf@check}[1]{\@ifundefined{qf@file@#1}{\PackageError{score2}{Unknown figure id #1}{}}{}}
\newcommand{\qf@img}[4][]{\includegraphics[width=#2\dimexpr\@nameuse{qf@w@#4}pt\relax,max width=#3,max height=0.45\textheight,#1]{\@nameuse{qf@file@#4}}}
% Full-width questions: finalize wraps every question that contains a wide figure in \qfullbegin...\qfullend,
% so the WHOLE question (stem, figures, options) is set across the page in reading order, then the two columns
% resume. The question's head - its number, stem and first figure (\qfullsplit marks the end of the head) -
% is kept on one page, so a number is never left at a page foot with its figure overleaf; the rest may break
% across pages normally, so no big gaps are left. Without a \qfullsplit, the whole question is kept together
% if it is under 45% of a page. Never write these yourself.
\newif\ifqf@full
\newif\ifqf@boxed
\newsavebox{\qf@box}
\newcommand{\qf@open}{\global\qf@fulltrue\par\global\qf@boxedtrue
  \setbox\qf@box\vbox\bgroup\hsize\textwidth\linewidth\textwidth}
\newcommand{\qf@emit}[1]{\par\egroup\global\qf@boxedfalse
  \ifdim\dimexpr\ht\qf@box+\dp\qf@box\relax<#1\textheight\noindent\usebox{\qf@box}\par\else\unvbox\qf@box\fi}
\newcommand{\qfullsplit}{\ifqf@boxed\qf@emit{0.85}\fi}
\newcommand{\qf@close}{\ifqf@boxed\qf@emit{0.45}\fi\global\qf@fullfalse}
\newcommand{\qfullbegin}{\end{multicols}\qf@open}
\newcommand{\qfullend}{\qf@close\begin{multicols}{2}}
\newcommand{\qfullstart}{\qf@open}
\newcommand{\qfulldone}{\qf@close}
% \qfig[scale]{id}: block figure, centred, directly where it is written (size set by the tool).
% Fallback if a wide figure is not inside a full-width question: break out of the columns in place - no floats.
\newcommand{\qfig}[2][1]{\qf@check{#2}%
  \ifqf@full\qf@blk{#1}{#2}\else
  \ifnum\@nameuse{qf@wide@#2}=1 \expandafter\@firstoftwo\else\expandafter\@secondoftwo\fi
  {\end{multicols}\begin{center}\qf@img{#1}{\textwidth}{#2}\end{center}\begin{multicols}{2}}%
  {\qf@blk{#1}{#2}}\fi}
\newcommand{\qf@blk}[2]{\par\smallskip{\centering\qf@img{#1}{\linewidth}{#2}\par}\smallskip}
% \qfigi[scale]{id}: inline figure (inside option lists / tables / minipages); never breaks columns.
\newcommand{\qfigi}[2][1]{\qf@check{#2}\qf@img[valign=c]{#1}{\linewidth}{#2}}
\makeatother
% structure macros - lint reads these to check completeness and figure placement
\newcommand{\papersec}[1]{\par\bigskip{\centering\Large\bfseries PAPER-#1\par}\medskip}
\newcommand{\partsec}[2]{\par\medskip\noindent{\large\bfseries PART-#1: #2}\par\smallskip}
\newcommand{\sect}[1]{\par\smallskip\noindent\textbf{#1}\par}
\newcommand{\qn}[1]{\par\smallskip\noindent\textbf{#1.}~}
\newcommand{\paraq}[2]{\par\smallskip\noindent\textbf{Paragraph for Questions #1 and #2}\par}
\newenvironment{qoptions}{\begin{itemize}[leftmargin=1.8em,itemsep=1pt,topsep=2pt,label={}]}{\end{itemize}}
\newcommand{\opt}[2]{\item[(#1)] #2}
% figure registry (generated by score2_tools.py finalize)
\qfigdef{p1_phy_s1_q01_1}{Score-II - Test 15 (09 Feb 2024)-img/p1_phy_s1_q01_1.png}{101.4}{0}
\qfigdef{p1_phy_s1_q01_2}{Score-II - Test 15 (09 Feb 2024)-img/p1_phy_s1_q01_2.png}{101.4}{0}
\qfigdef{p1_phy_s1_q01_3}{Score-II - Test 15 (09 Feb 2024)-img/p1_phy_s1_q01_3.png}{101.4}{0}
\qfigdef{p1_phy_s1_q01_4}{Score-II - Test 15 (09 Feb 2024)-img/p1_phy_s1_q01_4.png}{101.3}{0}
\qfigdef{p1_phy_s1_q06_1}{Score-II - Test 15 (09 Feb 2024)-img/p1_phy_s1_q06_1.png}{65.9}{0}
\qfigdef{p1_phy_s1_q06_2}{Score-II - Test 15 (09 Feb 2024)-img/p1_phy_s1_q06_2.png}{65.8}{0}
\qfigdef{p1_phy_s1_q06_3}{Score-II - Test 15 (09 Feb 2024)-img/p1_phy_s1_q06_3.png}{66.3}{0}
\qfigdef{p1_phy_s1_q06_4}{Score-II - Test 15 (09 Feb 2024)-img/p1_phy_s1_q06_4.png}{67.5}{0}
\qfigdef{p1_phy_s1_q09_1}{Score-II - Test 15 (09 Feb 2024)-img/p1_phy_s1_q09_1.png}{51.2}{0}
\qfigdef{p1_phy_s1_q09_2}{Score-II - Test 15 (09 Feb 2024)-img/p1_phy_s1_q09_2.png}{60.4}{0}
\qfigdef{p1_phy_s1_q09_3}{Score-II - Test 15 (09 Feb 2024)-img/p1_phy_s1_q09_3.png}{51.1}{0}
\qfigdef{p1_phy_s1_q09_4}{Score-II - Test 15 (09 Feb 2024)-img/p1_phy_s1_q09_4.png}{51.1}{0}
\qfigdef{p1_phy_s1_q11_1}{Score-II - Test 15 (09 Feb 2024)-img/p1_phy_s1_q11_1.png}{103.1}{0}
\qfigdef{p1_phy_s1_q12_1}{Score-II - Test 15 (09 Feb 2024)-img/p1_phy_s1_q12_1.png}{81.5}{0}
\qfigdef{p1_phy_s1_q12_2}{Score-II - Test 15 (09 Feb 2024)-img/p1_phy_s1_q12_2.png}{64.2}{0}
\qfigdef{p1_phy_s1_q12_3}{Score-II - Test 15 (09 Feb 2024)-img/p1_phy_s1_q12_3.png}{64.4}{0}
\qfigdef{p1_phy_s1_q12_5}{Score-II - Test 15 (09 Feb 2024)-img/p1_phy_s1_q12_5.png}{64.2}{0}
\qfigdef{p1_phy_s1_q12_4}{Score-II - Test 15 (09 Feb 2024)-img/p1_phy_s1_q12_4.png}{64.6}{0}
\qfigdef{p1_phy_s1_q14_1}{Score-II - Test 15 (09 Feb 2024)-img/p1_phy_s1_q14_1.png}{106.3}{0}
\qfigdef{p1_phy_s1_q17_1}{Score-II - Test 15 (09 Feb 2024)-img/p1_phy_s1_q17_1.png}{109.0}{0}
\qfigdef{p1_phy_s1_q19_1}{Score-II - Test 15 (09 Feb 2024)-img/p1_phy_s1_q19_1.png}{82.0}{0}
\qfigdef{p1_phy_s1_q20_1}{Score-II - Test 15 (09 Feb 2024)-img/p1_phy_s1_q20_1.png}{89.6}{0}
\qfigdef{p1_phy_s2_q04_1}{Score-II - Test 15 (09 Feb 2024)-img/p1_phy_s2_q04_1.png}{99.7}{0}
\qfigdef{p1_phy_s2_q05_1}{Score-II - Test 15 (09 Feb 2024)-img/p1_phy_s2_q05_1.png}{111.9}{0}
\qfigdef{p1_chem_s1_q09_1}{Score-II - Test 15 (09 Feb 2024)-img/p1_chem_s1_q09_1.png}{97.7}{0}
\qfigdef{p1_chem_s1_q09_2}{Score-II - Test 15 (09 Feb 2024)-img/p1_chem_s1_q09_2.png}{31.8}{0}
\qfigdef{p1_chem_s1_q09_3}{Score-II - Test 15 (09 Feb 2024)-img/p1_chem_s1_q09_3.png}{20.3}{0}
\qfigdef{p1_chem_s1_q09_4}{Score-II - Test 15 (09 Feb 2024)-img/p1_chem_s1_q09_4.png}{19.1}{0}
\qfigdef{p1_chem_s1_q09_5}{Score-II - Test 15 (09 Feb 2024)-img/p1_chem_s1_q09_5.png}{39.4}{0}
\qfigdef{p1_chem_s1_q10_1}{Score-II - Test 15 (09 Feb 2024)-img/p1_chem_s1_q10_1.png}{142.3}{0}
\qfigdef{p1_chem_s1_q10_A}{Score-II - Test 15 (09 Feb 2024)-img/p1_chem_s1_q10_A.png}{64.1}{0}
\qfigdef{p1_chem_s1_q10_B}{Score-II - Test 15 (09 Feb 2024)-img/p1_chem_s1_q10_B.png}{75.2}{0}
\qfigdef{p1_chem_s1_q10_C}{Score-II - Test 15 (09 Feb 2024)-img/p1_chem_s1_q10_C.png}{80.3}{0}
\qfigdef{p1_chem_s1_q10_D}{Score-II - Test 15 (09 Feb 2024)-img/p1_chem_s1_q10_D.png}{56.3}{0}
\qfigdef{p1_chem_s1_q11_1}{Score-II - Test 15 (09 Feb 2024)-img/p1_chem_s1_q11_1.png}{129.3}{0}
\qfigdef{p1_chem_s1_q11_3}{Score-II - Test 15 (09 Feb 2024)-img/p1_chem_s1_q11_3.png}{133.9}{0}
\qfigdef{p1_chem_s1_q11_4}{Score-II - Test 15 (09 Feb 2024)-img/p1_chem_s1_q11_4.png}{121.3}{0}
\qfigdef{p1_chem_s1_q11_5}{Score-II - Test 15 (09 Feb 2024)-img/p1_chem_s1_q11_5.png}{131.3}{0}
\qfigdef{p1_chem_s1_q14_1}{Score-II - Test 15 (09 Feb 2024)-img/p1_chem_s1_q14_1.png}{56.1}{0}
\qfigdef{p1_chem_s1_q14_2}{Score-II - Test 15 (09 Feb 2024)-img/p1_chem_s1_q14_2.png}{55.9}{0}
\qfigdef{p1_chem_s1_q14_3}{Score-II - Test 15 (09 Feb 2024)-img/p1_chem_s1_q14_3.png}{88.5}{0}
\qfigdef{p1_chem_s1_q14_4}{Score-II - Test 15 (09 Feb 2024)-img/p1_chem_s1_q14_4.png}{85.7}{0}
\qfigdef{p1_chem_s2_q05_1}{Score-II - Test 15 (09 Feb 2024)-img/p1_chem_s2_q05_1.png}{152.8}{0}
\qfigdef{p1_chem_s2_q05_2}{Score-II - Test 15 (09 Feb 2024)-img/p1_chem_s2_q05_2.png}{166.1}{0}
\qfigdef{p1_chem_s2_q05_3}{Score-II - Test 15 (09 Feb 2024)-img/p1_chem_s2_q05_3.png}{157.7}{0}
\qfigdef{p1_chem_s2_q10_1}{Score-II - Test 15 (09 Feb 2024)-img/p1_chem_s2_q10_1.png}{147.0}{0}
\qfigdef{p2_phy_s1_q04_1}{Score-II - Test 15 (09 Feb 2024)-img/p2_phy_s1_q04_1.png}{153.7}{0}
\qfigdef{p2_phy_s2_q05_1}{Score-II - Test 15 (09 Feb 2024)-img/p2_phy_s2_q05_1.png}{105.8}{0}
\qfigdef{p2_phy_s2_q06_1}{Score-II - Test 15 (09 Feb 2024)-img/p2_phy_s2_q06_1.png}{84.5}{0}
\qfigdef{p2_phy_s2_q09_1}{Score-II - Test 15 (09 Feb 2024)-img/p2_phy_s2_q09_1.png}{129.4}{0}
\qfigdef{p2_phy_s3_q03_1}{Score-II - Test 15 (09 Feb 2024)-img/p2_phy_s3_q03_1.png}{115.4}{0}
\qfigdef{p2_phy_s3_q04_1}{Score-II - Test 15 (09 Feb 2024)-img/p2_phy_s3_q04_1.png}{117.3}{0}
\qfigdef{p2_chem_s1_q04_1}{Score-II - Test 15 (09 Feb 2024)-img/p2_chem_s1_q04_1.png}{136.8}{0}
\qfigdef{p2_chem_s2_q09_1}{Score-II - Test 15 (09 Feb 2024)-img/p2_chem_s2_q09_1.png}{76.9}{0}
\qfigdef{p2_chem_s2_q09_2}{Score-II - Test 15 (09 Feb 2024)-img/p2_chem_s2_q09_2.png}{174.6}{0}
\qfigdef{p2_chem_s2_q09_4}{Score-II - Test 15 (09 Feb 2024)-img/p2_chem_s2_q09_4.png}{124.2}{0}
\qfigdef{p2_chem_s2_q09_3}{Score-II - Test 15 (09 Feb 2024)-img/p2_chem_s2_q09_3.png}{191.8}{0}
\qfigdef{p2_chem_s3_q02_1}{Score-II - Test 15 (09 Feb 2024)-img/p2_chem_s3_q02_1.png}{137.0}{0}
\qfigdef{p2_chem_s3_q06_1}{Score-II - Test 15 (09 Feb 2024)-img/p2_chem_s3_q06_1.png}{234.2}{0}
% ---- body skeleton ----
% \begin{document}
% \begin{center}{\Large\bfseries JEE (Main+Advanced) Score-II --- Test NN (DD Mon YYYY)}\end{center}
% \begin{multicols}{2}
%   \papersec{1}   (only for two-paper tests)
%   \partsec{1}{PHYSICS} \sect{Section-I (i): Single Correct} \qn{1} ... \qfig{p1_phy_q01_1} ...
% \end{multicols}
% \end{document}

\begin{document}
\begin{center}{\Large\bfseries JEE (Main+Advanced) Score-II --- Test 15 (09 Feb 2024)}\end{center}
\begin{multicols}{2}
\raggedcolumns
\papersec{1}
\partsec{1}{PHYSICS}
\sect{Section-I}

\qn{1} An electron beam is accelerated by a potential difference $V$ to hit a metallic target to produce X-rays. It produces continuous as well as characteristic X-rays. If $\lambda_{\min}$ is the smallest possible wavelength of X-ray in the spectrum, the variation of $\log\lambda_{\min}$ with $\log V$ is correctly represented in :-
\begin{qoptions}
\opt{A}{\qfigi{p1_phy_s1_q01_1}}
\opt{B}{\qfigi{p1_phy_s1_q01_2}}
\opt{C}{\qfigi{p1_phy_s1_q01_3}}
\opt{D}{\qfigi{p1_phy_s1_q01_4}}
\end{qoptions}

\qn{2} A nucleus of mass $M$ emits $\gamma$-ray photon of frequency `$\nu$'. The loss of internal energy by the nucleus is :\\
\relax[Take `$c$' as the speed of electromagnetic wave]
\begin{qoptions}
\opt{A}{$h\nu$}
\opt{B}{0}
\opt{C}{$h\nu\left[1-\dfrac{h\nu}{2Mc^2}\right]$}
\opt{D}{$h\nu\left[1+\dfrac{h\nu}{2Mc^2}\right]$}
\end{qoptions}

\qn{3} \textbf{Assertion (A) :} $\beta$-particles emitted from radioactive nuclei have continuous energy ranging from zero to a certain maximum value.\\
\textbf{Reason (R) :} In $\beta$-decay, a neutron is converted into a proton, an electron and an antineutrino. The total energy emitted in $\beta$-decay is shared by $\beta^-$ particle and antineutrino.
\begin{qoptions}
\opt{A}{(A) is true, (R) is true ; (R) is the correct explanation for (A).}
\opt{B}{(A) is true, (R) is true ; (R) is Not the correct explanation for (A).}
\opt{C}{(A) is True, (R) is False.}
\opt{D}{(A) is False, (R) is True.}
\end{qoptions}

\qn{4} When a hydrogen atom emits a photon of energy 12.1 eV, its orbital angular momentum changes by
\begin{qoptions}
\opt{A}{$1.05\times10^{-34}$ J s}
\opt{B}{$2.11\times10^{-34}$ J s}
\opt{C}{$3.16\times10^{-34}$ J s}
\opt{D}{$4.22\times10^{-34}$ J s}
\end{qoptions}

\qn{5} In the Bohr's model of hydrogen-like atom the force between the nucleus and the electron is modified as $F=\dfrac{e^2}{4\pi\epsilon_0}\left(\dfrac1{r^2}+\dfrac\beta{r^3}\right)$, where $\beta$ is a constant. For this atom, the radius of the $n^{\text{th}}$ orbit in terms of the Bohr radius $\left(a_0=\dfrac{\epsilon_0h^2}{m\pi e^2}\right)$ is :
\begin{qoptions}
\opt{A}{$r_n=a_0n-\beta$}
\opt{B}{$r_n=a_0n^2+\beta$}
\opt{C}{$r_n=a_0n+\beta$}
\opt{D}{$r_n=a_0n^2-\beta$}
\end{qoptions}

\qn{6} A long metallic bar is carrying heat from one of its ends to the other end under steady-state. The variation of temperature $\theta$ along the length $x$ of the bar from its hot end is best described by which of the following figures?
\begin{qoptions}
\opt{A}{\qfigi{p1_phy_s1_q06_1}}
\opt{B}{\qfigi{p1_phy_s1_q06_2}}
\opt{C}{\qfigi{p1_phy_s1_q06_3}}
\opt{D}{\qfigi{p1_phy_s1_q06_4}}
\end{qoptions}

\qn{7} One mole of an ideal gas (whose adiabatic constant is $\gamma$) undergoes the process $P=\alpha T^{1/2}$, where $\alpha$ is constant then find the molar specific heat of the gas.
\begin{qoptions}
\opt{A}{$\left(\dfrac{\gamma-1}{\gamma+1}\right)\dfrac R2$}
\opt{B}{$\left(\dfrac{\gamma-1}{\gamma+1}\right)^2R$}
\opt{C}{$\left(\dfrac{\gamma+1}{\gamma-1}\right)\dfrac R2$}
\opt{D}{$\left(\dfrac{\gamma+1}{\gamma}\right)R$}
\end{qoptions}

\qn{8} Hot water cools from 60$^\circ$C to 50$^\circ$C in the first 10 minutes and to 42$^\circ$C in the next 10 minutes. The temperature of the surroundings is :
\begin{qoptions}
\opt{A}{25$^\circ$C}
\opt{B}{20$^\circ$C}
\opt{C}{10$^\circ$C}
\opt{D}{15$^\circ$C}
\end{qoptions}

\qn{9} For a given amount of ideal gas :\par\nobreak
\begin{center}\small
\begin{tabular}{|c|c|c|p{2.7cm}|}\hline
 & \multicolumn{1}{c|}{\textbf{List-I}} & & \multicolumn{1}{c|}{\textbf{List-II}}\\\hline
(P) & \qfigi{p1_phy_s1_q09_1} & (1) & Heat is absorbed by system.\\\hline
(Q) & \qfigi{p1_phy_s1_q09_2} & (2) & Heat is rejected by system.\\\hline
(R) & \qfigi{p1_phy_s1_q09_3} & (3) & Internal energy will change.\\\hline
(S) & \qfigi{p1_phy_s1_q09_4} & (4) & Work is done by the system.\\\hline
\end{tabular}
\end{center}
\begin{qoptions}
\opt{A}{P$\to$1,4, Q$\to$3,4, R$\to$1,3,4, S$\to$2}
\opt{B}{P$\to$1,3,4, Q$\to$1,4, R$\to$1,4, S$\to$2}
\opt{C}{P$\to$3,4, Q$\to$1,3,4, R$\to$3,4, S$\to$1,3}
\opt{D}{P$\to$3,4, Q$\to$1,3,4, R$\to$3,4, S$\to$2}
\end{qoptions}

\qn{10} Consider a sample of oxygen behaving like an ideal gas. At 300 K, the ratio of root mean square (rms) velocity to the average velocity of gas molecule would be: (Molecular weight of oxygen is 32 g/mol; $R=8.3\ \mathrm{J\,K^{-1}\,mol^{-1}}$)
\begin{qoptions}
\opt{A}{$\sqrt{\dfrac33}$}
\opt{B}{$\sqrt{\dfrac83}$}
\opt{C}{$\sqrt{\dfrac{3\pi}8}$}
\opt{D}{$\sqrt{\dfrac{8\pi}3}$}
\end{qoptions}

\qn{11} A sample of gas follows process AB as shown in the figure. Then which of the following is incorrect about process AB ?
\qfig{p1_phy_s1_q11_1}
\begin{qoptions}
\opt{A}{$\Delta U=0$}
\opt{B}{$W>0$}
\opt{C}{$Q=W$}
\opt{D}{$\Delta U<0$}
\end{qoptions}

\qn{12} Figure shows a cyclic process. Which of the following is its PV conversion ?
\qfig{p1_phy_s1_q12_1}
\begin{qoptions}
\opt{A}{\qfigi{p1_phy_s1_q12_2}}
\opt{B}{\qfigi{p1_phy_s1_q12_3}}
\opt{C}{\qfigi{p1_phy_s1_q12_5}}
\opt{D}{\qfigi{p1_phy_s1_q12_4}}
\end{qoptions}

\qn{13} 200 g water is heated from 40$^\circ$C to 60$^\circ$C. Ignoring the slight expansion of water, the change in its internal energy is closed to (Given specific heat of water = 4184 J/kg/K) :
\begin{qoptions}
\opt{A}{4.2 kJ}
\opt{B}{8.4 kJ}
\opt{C}{167.4 kJ}
\opt{D}{16.7 kJ}
\end{qoptions}

\qn{14} P-V diagram of an ideal gas is as shown in figure. Work done by the gas in the process ABCD is :
\qfig{p1_phy_s1_q14_1}
\begin{qoptions}
\opt{A}{$4P_oV_o$}
\opt{B}{$2P_oV_o$}
\opt{C}{$3P_oV_o$}
\opt{D}{$P_oV_o$}
\end{qoptions}

\qn{15} \textbf{Statement-I :} If a gas container in motion is suddenly stopped, the temperature of the gas rises.\\
\textbf{Statement-II :} The kinetic energy of ordered mechanical motion is converted in to the kinetic energy of random motion of gas molecules.
\begin{qoptions}
\opt{A}{Statement$-$I is True, Statement$-$II is True ; Statement$-$II is a correct explanation for Statement$-$I.}
\opt{B}{Statement$-$I is True, Statement$-$II is True ; Statement$-$II is not a correct explanation for Statement$-$I.}
\opt{C}{Statement$-$I is True, Statement$-$II is False.}
\opt{D}{Statement$-$I is False, Statement$-$II is True.}
\end{qoptions}

\qn{16} Assuming the sun to be a spherical black body of radius $R$ at a temperature of $T$ K, evaluate the total radiant power, incident per unit area on earth, at a distance $r$ from the sun. $r_0$ is the radius of the earth and $\sigma$ is Stefan's constant.
\begin{qoptions}
\opt{A}{$\dfrac{R^2\sigma T^4}{r^2}$}
\opt{B}{$\dfrac{4\pi r_0^2R^2\sigma T^4}{r^2}$}
\opt{C}{$\dfrac{\pi r_0^2R^2\sigma T^4}{r^2}$}
\opt{D}{$\dfrac{R^2\sigma T^4}{4r^2}$}
\end{qoptions}

\qn{17} The spectrum of a black body at two temperature 27$^\circ$C and 327$^\circ$C is shown in the figure. Let $A_1$ and $A_2$ be the areas under the two curves respectively. The value of $\dfrac{A_2}{A_1}$ is :-
\qfig{p1_phy_s1_q17_1}
\begin{qoptions}
\opt{A}{1 : 16}
\opt{B}{4 : 1}
\opt{C}{2 : 1}
\opt{D}{16 : 1}
\end{qoptions}

\qn{18} Two identical masses of 5 kg each fall on a wheel from a height of 10m. The wheel disturbs a mass of 2 kg water, the rise in temperature of water will be :$-$
\begin{qoptions}
\opt{A}{2.6$^\circ$ C}
\opt{B}{1.2$^\circ$ C}
\opt{C}{0.32$^\circ$ C}
\opt{D}{0.12$^\circ$ C}
\end{qoptions}

\qn{19} Three immiscible liquids are filled in a container as shown. The base area of the container is $A$ and coefficient of cubical expansion of the material of the container is $\dfrac{3\gamma}2$ while the coefficient of cubical expansion of the liquids are shown in the figure. The temperature of the system is increased by $\Delta T$. The volume of the liquid flown out of the container is
\qfig{p1_phy_s1_q19_1}
\begin{qoptions}
\opt{A}{$\dfrac{A\ell\gamma\Delta T}3$}
\opt{B}{$A\lambda\gamma\Delta T$}
\opt{C}{$\dfrac{2A\ell\gamma\Delta T}3$}
\opt{D}{$\dfrac{A\ell\gamma\Delta T}2$}
\end{qoptions}

\qn{20} All the rods have same conductance `$K$' and same area of cross section `$A$'. If ends A and C are maintained at temperature $2T_0$ and $T_0$ respectively then which of the following is/are correct?
\qfig{p1_phy_s1_q20_1}
\begin{qoptions}
\opt{A}{Rate of heat flow through ABC, AOC and ADC is same}
\opt{B}{Rate of heat flow through BO and OD is not same}
\opt{C}{Total Rate of heat flow from A to C is $\dfrac{3KAT_0}{2a}$}
\opt{D}{Temperature at junctions B, O and D are same}
\end{qoptions}

\sect{Section-II}

\qn{1} Ice at $-20^\circ$C is filled upto height $h=10$ cm in a uniform cylindrical vessel. Water at temperature $\theta^\circ$C is filled in another identical vessel upto the same height $h=10$ cm. Now, water from second vessel is poured into first vessel and it is found that level of upper surface falls through $\Delta h=0.5$ cm when thermal equilibrium is reached. Neglecting thermal capacity of vessels, change in density of water due to change in temperature and loss of heat due to radiation, calculate initial temperature $\theta$ of water. (in $^\circ$C)\\
Given :\\
Density of water: $\rho_w=1\ \mathrm{gm\,cm^{-3}}$\\
Density of ice: $\rho_i=0.9\ \mathrm{gm/cm^3}$\\
Specific heat of water: $s_w=1$ cal/gm $^\circ$C\\
Specific heat of ice: $s_i=0.5$ cal/gm$^\circ$C\\
Specific latent heat of ice: $L=80$ cal/gm

\qn{2} The volume of a certain mass of gas at constant pressure is doubled to its value at 0$^\circ$C. The temperature (in $^\circ$C) of gas will be :

\qn{3} A container is divided into two chambers by a partition. The volume of first chamber is 4.5 litre and second chamber is 5.5 litre. The first chamber contain 3.0 moles of gas at pressure 2.0 atm and second chamber contain 4.0 moles of gas at pressure 3.0 atm. After the partition is removed and the mixture attains equilibrium, then, the common equilibrium pressure existing in the mixture is $x\times10^{-1}$ atm. Value of $x$ is\underline{\hspace{2em}}.

\qn{4} When heat is supplied to the gas it expands and displaces piston by $L/2$ where natural length of springs are $L=1$ m. Spring constant $K=100$ N/m. Area of piston is 1 m$^2$. The pressure of gas in final situation is (in N/m$^2$) -
\qfig{p1_phy_s2_q04_1}

\qn{5} For the given cyclic process CAB as shown for a gas, the work done (in J) is :
\qfig{p1_phy_s2_q05_1}

\qn{6} If $2.5\times10^{-6}$ N average force is exerted by a light wave on a non-reflecting surface of 30 cm$^2$ area during 40 minutes of time span, the energy flux of light just before it falls on the surface is \underline{\hspace{2em}} W/cm$^2$.\\
(Round off to the Nearest Integer) (Assume complete absorption and normal incidence conditions are there)

\qn{7} A He$^+$ ion in ground state is fired towards a Hydrogen atom in ground state and at rest. What should be the minimum kinetic energy (in eV) of He$^+$ ion so that both single electron species may get excited?

\qn{8} When the voltage applied to an x-ray tube is increased from $V_1=10$kV to $V_2=20$ kV, the wavelength difference of K$_\alpha$ line and cut-off wavelength of continuous x-ray spectrum increases by a factor $\eta$. If atomic number of the element of which the tube's anode is made, is 29 then find the value of $\eta$.

\qn{9} From the given data, the amount of energy required to break the nucleus of aluminium is \underline{\hspace{2em}} $\times10^{-3}$ J.\\
Mass of neutron = 1.00866 u; Mass of proton = 1.00726 u; Mass of Aluminium nucleus = 27.18846 u;\\
(Assume 1 u corresponds to $x$ J of energy)\\
(Round off to the nearest integer)

\qn{10} The respective speeds of five molecules are 2, 1.5, 1.6, 1.6 and 1.2 km/sec. The most probable speed is $x$ km/sec, then $10x$ will be :-

\partsec{2}{CHEMISTRY}
\sect{Section-I}

\qn{1} At the same temperature and pressure, oxygen gas is more soluble in water than hydrogen gas. Select the correct statement.
\begin{qoptions}
\opt{A}{Henry constant (K$_\mathrm{H}$) depend on temperature only.}
\opt{B}{Henry constant (K$_\mathrm{H}$) does not depend on nature of gas.}
\opt{C}{Henry constant (K$_\mathrm{H}$) of oxygen is greater than hydrogen.}
\opt{D}{Henry constant (K$_\mathrm{H}$) of oxygen is lesser than hydrogen.}
\end{qoptions}

\qn{2} Which of the following solution show higher observed boiling point then expected boiling point :
\begin{qoptions}
\opt{A}{Solution of acetone and chloroform}
\opt{B}{Solution of carbon tetrachloride and chloroform}
\opt{C}{Solution of acetone and benzene}
\opt{D}{Solution of acetone and carbon-di-sulphide}
\end{qoptions}

\qn{3} Mole fraction of component B in vapour phase is $y_1$ and that of component A in liquid phase is $x_1$ the the total vapour pressure of liquid mixture is :\\
($\mathrm{P^\circ_A}$ = vapour pressure of A and $\mathrm{P^\circ_B}$ = vapour pressure of B)
\begin{qoptions}
\opt{A}{$\dfrac{x_1P_A^0}{y_1}$}
\opt{B}{$\dfrac{y_1P_B^0}{x_1}$}
\opt{C}{$\dfrac{x_1P_A^0}{1-y_1}$}
\opt{D}{$\dfrac{x_1P_B^0}{y_1}$}
\end{qoptions}

\qn{4} The relative lowering of the vapour pressure of an aqueous solution containing a urea solute is 0.2. The molality of solution is :
\begin{qoptions}
\opt{A}{$\dfrac{1000}{72}$ mole/kg}
\opt{B}{$\dfrac{1000}9$ mole/kg}
\opt{C}{$\dfrac2{10}$ mole/kg}
\opt{D}{$\dfrac{100}3$ mole/kg}
\end{qoptions}

\qn{5} 1 molal aqueous solution of trichloroacetic acid is heated to its boiling point. The solution has the boiling point of 100.65$^\circ$C. The dissociation constant of trichloroacetic acid is :\\
(Given : K$_\mathrm{b}$ of water is 0.52 K-kg/mol,\\
molality = molarity of solution)
\begin{qoptions}
\opt{A}{$1\times10^{-5}$}
\opt{B}{$8.33\times10^{-2}$}
\opt{C}{$6.25\times10^{-2}$}
\opt{D}{$3.33\times10^{-2}$}
\end{qoptions}

\qn{6} Deci-molal aqueous solution of potassium ferrocynaide is 50\% dissociated. What will be its freezing point at 1 atm.\\
(Given : K$_\mathrm{f}$ of water is 1.86 K-kg/mol)
\begin{qoptions}
\opt{A}{0.558 K}
\opt{B}{0.558$^\circ$C}
\opt{C}{$-0.186$ K}
\opt{D}{$-0.558^\circ$C}
\end{qoptions}

\qn{7} \textbf{Statement-1 :} 5\%w/V aqueous solution of glucose is isotonic with 5\%w/V aqueous solution of urea at same temperature.\\
\textbf{Statement-2 :} Glucose and urea has same Van't Hoff factor in their own solution.
\begin{qoptions}
\opt{A}{Statement-1 is true, Statement-2 is true and Statement-2 is correct explanation for Statement-1}
\opt{B}{Statement-1 is true, Statement-2 is true and Statement-2 is NOT the correct explanation for Statement-1}
\opt{C}{Statement-1 is true, Statement-2 is false}
\opt{D}{Statement-1 is false, Statement-2 is true}
\end{qoptions}

\qn{8} Which one of the following set of elements can be detected using sodium fusion extract ?
\begin{qoptions}
\opt{A}{Sulfur, Nitrogen, Phosphorous, Halogens}
\opt{B}{Phosphorous, Oxygen, Nitrogen, Halogens}
\opt{C}{Nitrogen, Phosphorous, Carbon, Sulfur}
\opt{D}{Halogens, Nitrogen, Oxygen, Sulfur}
\end{qoptions}

\qn{9}\par\nobreak
\qfig{p1_chem_s1_q09_1}
\begin{qoptions}
\opt{A}{Compound (B) is\\ \qfigi{p1_chem_s1_q09_2}}
\opt{B}{Compound (B) is\\ \qfigi{p1_chem_s1_q09_3}}
\opt{C}{Compound (A) is\\ \qfigi{p1_chem_s1_q09_4}}
\opt{D}{Compound (B) is\\ \qfigi{p1_chem_s1_q09_5}}
\end{qoptions}

\qn{10} Identify the major products A and B respectively in the following reactions of phenol.
\qfig{p1_chem_s1_q10_1}
\begin{qoptions}
\opt{A}{\qfigi{p1_chem_s1_q10_A}}
\opt{B}{\qfigi{p1_chem_s1_q10_B}}
\opt{C}{\qfigi{p1_chem_s1_q10_C}}
\opt{D}{\qfigi{p1_chem_s1_q10_D}}
\end{qoptions}

\qn{11} Out of the following which is correct :
\begin{qoptions}
\opt{A}{\qfigi{p1_chem_s1_q11_1}}
\opt{B}{\qfigi{p1_chem_s1_q11_3}}
\opt{C}{\qfigi{p1_chem_s1_q11_4}}
\opt{D}{\qfigi{p1_chem_s1_q11_5}}
\end{qoptions}

\qn{12} Thermosetting polymer, Bakelite is formed by the reaction of phenol with :-
\begin{qoptions}
\opt{A}{CH$_3$CH$_2$CHO}
\opt{B}{CH$_3$CHO}
\opt{C}{HCHO}
\opt{D}{HCOOH}
\end{qoptions}

\qn{13} Which of the following is correct about H-bonding in nucleotide ?
\begin{qoptions}
\opt{A}{A----A and T---T}
\opt{B}{G----T and A----C}
\opt{C}{A --- G and T ----C}
\opt{D}{A --- T and G ---- C}
\end{qoptions}

\qn{14}\par\nobreak
\begin{center}\small
\begin{tabular}{|c|p{2.6cm}|c|p{2.5cm}|}\hline
 & \multicolumn{1}{c|}{\textbf{List-I}} & & \multicolumn{1}{c|}{\textbf{List-II}}\\\hline
(P) & \qfigi{p1_chem_s1_q14_1} & (1) & Product give molish test\\\hline
(Q) & \qfigi{p1_chem_s1_q14_2} & (2) & Product give violet colour with neutral FeCl$_3$\\\hline
(R) & \qfigi{p1_chem_s1_q14_3}\newline Starch product & (3) & L.S. of product give blue colour solution with FeSO$_4$ solution\\\hline
(S) & \qfigi{p1_chem_s1_q14_4} & (4) & Product give Biurate test\\\hline
\end{tabular}
\end{center}
\begin{qoptions}
\opt{A}{P $\to$ 3;Q $\to$ 2,3;R $\to$ 1;S $\to$ 4}
\opt{B}{P $\to$ 2;Q $\to$ 2,3;R $\to$ 1;S $\to$ 3,4}
\opt{C}{P $\to$ 3;Q $\to$ 2,3;R $\to$ 1;S $\to$ 3,4}
\opt{D}{P $\to$ 3;Q $\to$ 2;R $\to$ 4;S $\to$ 1,3}
\end{qoptions}

\qn{15} Which of the following is correct order of solubility of given salts in aq solution?
\begin{qoptions}
\opt{A}{Be(NO$_3$)$_2$ $<$ Mg(NO$_3$)$_2$ $<$ Ca(NO$_3$)$_2$ $<$ Sr(NO$_3$)$_2$ $<$ Ba(NO$_3$)$_2$}
\opt{B}{Li$_2$CO$_3$ $<$ Na$_2$CO$_3$ $<$ K$_2$CO$_3$ $<$ Rb$_2$CO$_3$ $<$ Cs$_2$CO$_3$}
\opt{C}{LiF $>$ NaF $>$ KF $>$ RbF $>$ CsF}
\opt{D}{AgF $<$ AgCl $<$ AgBr $<$ AgI}
\end{qoptions}

\qn{16} Which statement is correct among the following?
\begin{qoptions}
\opt{A}{BF$_3$ in strongest Lewis acid among BF$_3$, BCl$_3$, BBr$_3$}
\opt{B}{Between NH$_3$ and PH$_3$, PH$_3$ is a better electron donor because PH$_3$ has more directional lone pair}
\opt{C}{C$\ell$OF$_3$ has same shape as XeF$_6$}
\opt{D}{The magnitude of repulsion between bonding pairs of electrons depends on the electronegativity difference between the central atom and other atoms}
\end{qoptions}

\qn{17} Select the correct statement among the following options :-
\begin{qoptions}
\opt{A}{Type of overlapping in the bridge bond is same in Al$_2$(CH$_3$)$_6$ , Al$_2$ Cl$_6$ , Be$_2$Cl$_4$}
\opt{B}{Back bonding is present in H$_3$SiNCS, BF$_3$ , PF$_3$ and N(SiH$_3$)$_3$}
\opt{C}{I$_2$Cl$_6$ is a planar molecule having maximum 6 atoms in a plane}
\opt{D}{Hybridisation of central atom in C$\ell$O$_2$ is sp$^3$}
\end{qoptions}

\qn{18} Which of the following molecular orbitals has two nodal plane?
\begin{qoptions}
\opt{A}{$\pi_{2\mathrm{P_y}}$}
\opt{B}{$\pi^*_{2\mathrm{P_x}}$}
\opt{C}{$\sigma^*_{2\mathrm{P_z}}$}
\opt{D}{3d$_{xy}$}
\end{qoptions}

\qn{19} \textbf{Statement-1 :-} Both PCl$_5$ and BrF$_5$ have identical shapes.\\
\textbf{Statement-2 :-} Same number of $\sigma$ bonds are present in both PCl$_5$ and BrF$_5$
\begin{qoptions}
\opt{A}{Statement-1 is true, statement-2 is true and statement-2 is correct explanation for statement-1.}
\opt{B}{Statement-1 is true, statement-2 is true and statement-2 is NOT the correct explanation for statement-1.}
\opt{C}{Statement-1 is true, statement-2 is false.}
\opt{D}{Statement-1 is false, statement-2 is true.}
\end{qoptions}

\qn{20} Which of the following order of compound is correct against to their indicated property :-
\begin{qoptions}
\opt{A}{PbCl$_2$ $<$ PbCl$_4$ (Stability)}
\opt{B}{FeCl$_2$ $>$ FeCl$_3$ (Covalent nature)}
\opt{C}{LiF $>$ NaF (Thermal stability)}
\opt{D}{N$_2$O $>$ NO (Acidic nature)}
\end{qoptions}

\sect{Section-II}

\qn{1} Find the value of $Z$ for 19.6\%w/V of sulphuric acid having density 1.096.\\
$Z$ = relative lowering in vapour pressure $\times10^4$

\qn{2} A solution prepared by mixing 92 gram glycerin (C$_3$H$_8$O$_3$) to 408 gram of water. Solution has density 0.8 g/ml at 300 K. (Given : $R=0.08$ atm-L/mole-K)\\
Find $Z$.\\
\relax[$Z$ = osmotic pressure (in atm) $\times10$]

\qn{3} 200 g of acetic acid was contaminated with 17.6 g ascorbic acid (C$_6$H$_8$O$_6$) to lower down its freezing point by ($z\times10^{-3}$)$^\circ$C. The value of $z$ is\underline{\hspace{3em}}.\\
(K$_f$ = 3.9 K-kg/mol)

\qn{4} Identify total number of Essential amino acids\\
(a) Glycine \quad (b) Alanine \quad (c) Phenyl alanine\\
(d) Arginine \quad (e) Cysteine \quad (f) Histidine\\
(g) Proline \quad (h) Tryptophan

\qn{5} Number of compounds which can show faster rate of nucleophilic substitution reaction with NaOCH$_3$
\qfig{p1_chem_s2_q05_1}
\qfig{p1_chem_s2_q05_2}
\qfig{p1_chem_s2_q05_3}

\qn{6} How many of the following have been correctly matched with their property.\\
(a) ClO$_2$ $-$ Contains odd electron but does not undergo dimerization.\\
(b) XeF$_8^{-2}$ $-$ Does not exist.\\
(c) Cl$_2$O $-$ Bond angle greater than 109$^\circ$28.\\
(d) B$_2$ $-$ Diamagnetic if s-p mixing in not considered.\\
(e) NO $-$ Diamagnetic gas\\
(f) N$_2$H$_4$ - N-N bond length greater than N$_2$F$_4$.\\
(g) H$_2$S$_2$O$_8$ $-$ Contains peroxy linkage.\\
(h) XeF$_4$ $-$ d$_{x^2-y^2}$ orbital is involved in hybridisation of central atom.

\qn{7} Find the number of molecules which is/are planar and having dipole moment\\
SF$_4$ , PCl$_3$F$_2$ , PCl$_2$F$_3$ , I$_3^-$ , BrF$_3$ , ICl$_2^-$ , XeF$_4$ , O$_3$ , ClO$_2$ , NO$_2^-$, NO$_3^-$.

\qn{8} Ag$_2$CO$_3$ on thermal decomposition produces two gases A and B. Gas B is paramagnetic in nature. If the total number of e$^-$ present in all ABMO in divalent anion of gas B is $X$ and the total number of lone pair present on gas A is $Y$, then :-\\
Find the value of $\dfrac{(2X-Y)}2$

\qn{9} Number of molecules or ions which has/have any bond angle of 120$^\circ$ is :\\
$\overset{-}{\mathrm{C}}$H$_3$, $\overset{\oplus}{\mathrm{C}}$H$_3$, NH$_3$, CO$_3^{2-}$, BCl$_3$, PCl$_5$

\qn{10} How many of the following groups are capable of donating electron density to the benzene ring during the Wheland intermediate, being ortho / para directors :\\
(a) $-$Cl \hfill (b) $-$NO \hfill (c) NO$_2$\\
(d) $-$CCl$_3$ \hfill (e) $-$CH=CH$_2$\\
(f) $-$CH = CH $-$ Cl. \hfill (g) $-$CH = CH $-$ NO$_2$
\qfig{p1_chem_s2_q10_1}

\partsec{3}{MATHEMATICS}
\sect{Section-I}

\qn{1} If $A$ is square matrix of order 3 such that $|A|=3$, then $|\text{adj}\,(\text{adj}\,(\text{adj}\,A))|$ is-
\begin{qoptions}
\opt{A}{$3^8$}
\opt{B}{$3^9$}
\opt{C}{$3^{10}$}
\opt{D}{$3^{11}$}
\end{qoptions}

\qn{2} If $a,b,c,d,\alpha,\beta>0$, such that $\dfrac ab=\dfrac cd=\lambda$, then $\lambda$ is not always equal to
\begin{qoptions}
\opt{A}{$\sqrt{\dfrac{(\alpha a)^2+(\beta c)^2}{(\alpha b)^2+(\beta d)^2}}$}
\opt{B}{$\left(\dfrac{a^4+c^4}{b^4+d^4}\right)^{1/4}$}
\opt{C}{$\left(\dfrac{\alpha a+\beta c}{\alpha b+\beta d}\right)$}
\opt{D}{$\sqrt{\dfrac{ad+bc}{ac+bd}}$}
\end{qoptions}

\qn{3} Let $a,b,c$ be three real numbers which satisfy $a^2+b^2+c^2+ab+bc+ca=0$.\\
\textbf{Assertion (A) :} $\begin{vmatrix}a+b&a+2&b-3\\b-2&b+c&c-2\\c+3&a+2&c+a\end{vmatrix}=0$.\\
\textbf{and}\\
\textbf{Reason (R) :} The value of skew symmetric determinant is always zero.
\begin{qoptions}
\opt{A}{Both \textbf{A} and \textbf{R} are correct and \textbf{R} is the correct explanation of \textbf{A}}
\opt{B}{Both \textbf{A} and \textbf{R} are correct but \textbf{R} is \textbf{NOT} the correct explanation of \textbf{A}}
\opt{C}{\textbf{A} is correct but \textbf{R} is not correct}
\opt{D}{\textbf{A} is not correct but \textbf{R} is correct}
\end{qoptions}

\qn{4} If $A$ is a square matrix such that $A^2=A$, then $(I+A)^3-7A$ is\\
(where $I$ is an identity matrix of same order as of $A$)
\begin{qoptions}
\opt{A}{$3I$}
\opt{B}{$O$}
\opt{C}{$I$}
\opt{D}{$2I$}
\end{qoptions}

\qn{5} If $1,\alpha_1,\alpha_2,\alpha_3,\dots,\alpha_8$ are nine, ninth roots of unity (taken in counter clockwise sequence) then $\left|(2-\alpha_1)(2-\alpha_3)(2-\alpha_5)(2-\alpha_7)\right|$ is equal to
\begin{qoptions}
\opt{A}{$\sqrt{255}$}
\opt{B}{$\sqrt{511}$}
\opt{C}{$\sqrt{1023}$}
\opt{D}{15}
\end{qoptions}

\qn{6} In $\Delta$ABC, $y=x$ is the internal angle bisector of angle B and equation of side AC is $2x-y=2$. If coordinates of A are (4, 6) and $2\text{AB}=\text{BC}$, then coordinates of B is
\begin{qoptions}
\opt{A}{(4, 4)}
\opt{B}{(14, 14)}
\opt{C}{$\left(\dfrac{14}9,\dfrac{14}9\right)$}
\opt{D}{$\left(\dfrac49,\dfrac49\right)$}
\end{qoptions}

\qn{7} From point P(5, 1) a line is drawn to cut the circle $x^2+y^2+4y-21=0$ at A and B such that PA = 1, then the length of perpendicular from centre of circle upon chord AB, is
\begin{qoptions}
\opt{A}{3}
\opt{B}{4}
\opt{C}{$\sqrt6$}
\opt{D}{2}
\end{qoptions}

\qn{8} If $\begin{vmatrix}1&x&x^2\\x&x^2&1\\x^2&1&x\end{vmatrix}=3$ then the value of $\begin{vmatrix}x^3-1&0&x-x^4\\0&x-x^4&x^3-1\\x-x^4&x^3-1&0\end{vmatrix}$ is
\begin{qoptions}
\opt{A}{0}
\opt{B}{3}
\opt{C}{6}
\opt{D}{9}
\end{qoptions}

\qn{9}\par\nobreak
\begin{center}\small
\begin{tabular}{|c|p{3.6cm}|c|c|}\hline
 & \multicolumn{1}{c|}{\textbf{List-I}} & & \textbf{List-II}\\\hline
(A) & If $z$ satisfy $|z+3-i|=4$ and $\arg(z)=\dfrac{5\pi}4$, then $z$ is & (I) & $1-i$\\\hline
(B) & $z$ is a complex number that satisfy $z+2\bar z=3+i$, then $z$ is & (II) & $-3-3i$\\\hline
(C) & $\dfrac{(1+i)^{8n+5}}{(1-i)^{8n+4}}$ is equal to & (III) & $3-\dfrac i2$\\\hline
(D) & If $z$ satisfy $\arg\left(\dfrac{\sqrt2-z}{\sqrt2+z}\right)=\pm\dfrac\pi2$, then $z$ can be & (IV) & $1+i$\\\hline
\end{tabular}
\end{center}
\begin{qoptions}
\opt{A}{(A)$\to$(II), (B)$\to$(I), (C)$\to$(IV), (D)$\to$(I,IV)}
\opt{B}{(A)$\to$(I,IV), (B)$\to$(II), (C)$\to$(I), (D)$\to$(IV)}
\opt{C}{(A)$\to$(IV), (B)$\to$(I,IV), (C)$\to$(II), (D)$\to$(I)}
\opt{D}{(A)$\to$(I), (B)$\to$(IV), (C)$\to$(I,IV), (D)$\to$(II)}
\end{qoptions}

\qn{10} Let $|z_1|=1$, $|z_2|=2$, $|z_3|=3$ and $z_1+z_2+z_3=3+\sqrt5\,i$ then the value of $\text{Re}\left(z_1\bar z_2+z_2\bar z_3+z_3\bar z_1\right)$ is equal to\\
(where $z_1,z_2$ and $z_3$ are complex numbers and $i=\sqrt{-1}$)
\begin{qoptions}
\opt{A}{1}
\opt{B}{$-1$}
\opt{C}{$\dfrac{-1}2$}
\opt{D}{0}
\end{qoptions}

\qn{11} If the point P$\left(\dfrac{3\lambda}2,1\right)$ lies between the two different lines $x+y=\lambda$ and $x+y=2\lambda$ then the least positive integral value of $\lambda$, is
\begin{qoptions}
\opt{A}{3}
\opt{B}{4}
\opt{C}{0}
\opt{D}{5}
\end{qoptions}

\qn{12} Number of point of intersection of the curves $\arg(z-3i)=\dfrac{3\pi}4$ and $\arg(2z+1-2i)=\dfrac\pi2$, is(are)
\begin{qoptions}
\opt{A}{0}
\opt{B}{1}
\opt{C}{2}
\opt{D}{3}
\end{qoptions}

\qn{13} If $a,b,c$ are real then the value of determinant $\begin{vmatrix}a^2+1&ab&ac\\ab&b^2+1&bc\\ac&bc&c^2+1\end{vmatrix}=1$ if
\begin{qoptions}
\opt{A}{$a+b+c=2$}
\opt{B}{$a+b+c=1$}
\opt{C}{$a+b+c=-1$}
\opt{D}{$a=b=c=0$}
\end{qoptions}

\qn{14} If $z_1$ and $z_2$ both satisfy the equation $\text{Re}(z)=|z-2|$ and $\arg(z_1-z_2)=\dfrac\pi3$ then $\text{Im}(z_1+z_2)$ is\\
(where $\text{Im}(z)$ represents imaginary part of $z$)
\begin{qoptions}
\opt{A}{$\dfrac1{\sqrt3}$}
\opt{B}{$\dfrac2{\sqrt3}$}
\opt{C}{$\sqrt3$}
\opt{D}{$\dfrac4{\sqrt3}$}
\end{qoptions}

\qn{15} Let $A=\begin{bmatrix}0&-3\\0&0\end{bmatrix}$ and $I$ be unit matrix of order 2.\\
\textbf{Statement-1 :} $(I+A)^{532}=I+532A$\\
\textbf{and}\\
\textbf{Statement-2 :} If $B$ be a nilpotent matrix of index 2 and $I$ be a unit matrix of same order as $B$ then $(I+B)^n=I+nB$, $n\in N$.
\begin{qoptions}
\opt{A}{Statement-1 is True, Statement-2 is True}
\opt{B}{Statement-1 is False, Statement-2 is False}
\opt{C}{Statement-1 is True, Statement-2 is False.}
\opt{D}{Statement-1 is False, Statement-2 is True.}
\end{qoptions}

\qn{16} Let $z$ be a complex number satisfying $|z+16|=4|z+1|$. Then
\begin{qoptions}
\opt{A}{$|z|=4$}
\opt{B}{$|z|=5$}
\opt{C}{$|z|=6$}
\opt{D}{$|z|=7$}
\end{qoptions}

\qn{17} Consider a $3\times3$ matrix $A=\begin{bmatrix}1&0&0\\0&\alpha&\beta\\0&\beta&\alpha\end{bmatrix}$ and $\left|2A^3\right|=2^{15}$ ($\alpha,\beta\in I$), then maximum value of $\alpha$, is
\begin{qoptions}
\opt{A}{2}
\opt{B}{3}
\opt{C}{4}
\opt{D}{5}
\end{qoptions}

\qn{18} If $w$ is the cube root of unity and two of the factors for $\begin{vmatrix}a&b&c\\b&c&a\\c&a&b\end{vmatrix}$ are $a+bw^p+cw^q$ and $a+bw^q+cw^p$ then $p+q=$
\begin{qoptions}
\opt{A}{1}
\opt{B}{2}
\opt{C}{3}
\opt{D}{0}
\end{qoptions}

\qn{19} The circle $S\equiv x^2+y^2+kx+(k+1)y-(k+1)=0$ always passes through two fixed point for every real $k$. If the minimum value of the radius of the circle $S$ is $\dfrac1{\sqrt P}$ then the value of `P' is :
\begin{qoptions}
\opt{A}{4}
\opt{B}{6}
\opt{C}{8}
\opt{D}{16}
\end{qoptions}

\qn{20} The smallest radius of circle which orthogonally intersects the circles $x^2+y^2=4$ and $(x-6)^2+y^2=2$, is
\begin{qoptions}
\opt{A}{$\dfrac{\sqrt{119}}6$}
\opt{B}{$\dfrac{\sqrt{217}}6$}
\opt{C}{$\dfrac{\sqrt{223}}6$}
\opt{D}{$\dfrac{\sqrt{131}}6$}
\end{qoptions}

\sect{Section-II}

\qn{1} If values of $\lambda$ and $\mu$ are such that the distance `$d$' between the parallel lines $L_1$ and $L_2$ is maximum, where\\
$L_1:(x+y-3)+\lambda(3x-2y+1)=0$\\
$L_2:(x-2y-8)+\mu(2x-y-10)=0$\\
then value of $d$ is equal to

\qn{2} Let $A$ and $B$ be two $2\times2$ non singular skew symmetric matrices such that $AB=BA$. Where $P^T$ denotes transpose of $P$ if $\left(B^2A^2\right)^2\left(\left(B^TA\right)^{-1}\right)^2\left(\left(BA^{-1}\right)^T\right)^2=B^n$ then $n$ is equal to

\qn{3} If $z_1=8+2i$ and $z_2=4+2i$ are two complex numbers on the complex plane. A complex number $z$ satisfying $\text{amp}\left(\dfrac{z-z_1}{z-z_2}\right)=\dfrac\pi6$ moves on the major arc of a circle whose radius, is

\qn{4} Let $M$ is a $3\times3$ matrix such that $|M|=2$ and if $\text{adj}\left(\text{adj}\left(|M|M^2\right)\right)=2^\alpha M^\beta$ then $(\alpha-2\beta)$ is equal to

\qn{5} Consider a triangle ABC with A(0, 0), B($2\sqrt3$, 0) and C($\sqrt3$, 3). If L is locus of the points from which tangents of equal lengths can be drawn to its incircle \& excircle opposite to vertex A, then the minimum distance of L from origin, is

\qn{6} A polynomial $P(x)$ of degree 5 is such that $P(0)=-120$, $P(1)=1$, $P(2)=16$, $P(3)=81$, $P(4)=256$ and $P(5)=625$, then the value of $P'(3)$ is (where $P'(x)$ represents differentiation of $P(x)$)

\qn{7} If $A=\begin{bmatrix}1&1&1\\0&1&1\\0&0&1\end{bmatrix}$ and $M=A+A^2+A^3+\dots..+A^{20}$, then the sum of all the elements of the matrix $M$ is equal to \underline{\hspace{3em}}

\qn{8} Two side of parallelogram are along the lines $4x+5y=0$ and $7x+2y=0$. If the equation of one of the diagonal is $11x+7y=9$, then other diagonal having equation $y=mx$ where $m$ is equal to

\qn{9} In a $\Delta$ABC, the vertex A is (1, 1) and is orthocenter is (2, 4). If the sides AB and BC are members of the family of straight lines $ax+by+c=0$, where $a,b,c$ are in A.P., then the coordinates of vertex C are $(h,k)$. The value of $2h+12k$ is equal to \underline{\hspace{3em}}.

\qn{10} Area of the parallelogram having sides $14x-21y+15=0$, $14x-21y-6=0$, $2x-6y+9=0$ \& $2x-6y-3=0$, is equal to

\papersec{2}
\partsec{1}{PHYSICS}
\sect{Section-I(i)}

\qn{1} A particle A of mass `m' and charge `q' is accelerated by a potential difference of 50 V. Another particle B of mass `4 m' and charge `q' is accelerated by a potential difference of 2500 V. The ratio of de-Broglie wavelengths $\dfrac{\lambda_A}{\lambda_B}$ is close to :
\begin{qoptions}
\opt{A}{10.00}
\opt{B}{14.14}
\opt{C}{4.47}
\opt{D}{0.07}
\end{qoptions}

\qn{2} 50 W/m$^2$ energy density of sunlight is normally incident on the surface of a solar panel. Some part of incident energy (25\%) is reflected from the surface and the rest is absorbed. The force exerted on 1m$^2$ surface area will be close to ($c = 3\times10^8$ m/s) :-
\begin{qoptions}
\opt{A}{$15\times10^{-8}$ N}
\opt{B}{$35\times10^{-8}$ N}
\opt{C}{$10\times10^{-8}$ N}
\opt{D}{$20\times10^{-8}$ N}
\end{qoptions}

\qn{3} An electron of a hydrogen like atom, having $Z = 4$, jumps from 4$^{\text{th}}$ energy state to 2$^{\text{nd}}$ energy state. The energy released in this process, will be: (Given $Rch = 13.6$ eV); where $R$ = Rydberg constant, $c$ = Speed of light in vacuum, $h$ = Planck's constant
\begin{qoptions}
\opt{A}{13.6 eV}
\opt{B}{10.5 eV}
\opt{C}{3.4 eV}
\opt{D}{40.8 eV}
\end{qoptions}

\qn{4} A gas fills the right portion of a horizontal cylinder whose radius is 5.00 cm. The initial pressure of the gas is 101 kPa. A frictionless movable piston separates the gas from the left portion of the cylinder, which is evacuated and contains an ideal spring, as the drawing shows. The piston is initially held in place by a pin. The spring is initially unstrained, and the length of the gas-filled portion is 20.0 cm. When the pin is removed and the gas is allowed to expand, the length of the gas-filled chamber doubles when it finally reaches equilibrium. The initial and final temperatures are equal. Determine the spring constant of the spring (nearly).
\par\nobreak
\qfig{p2_phy_s1_q04_1}
\begin{qoptions}
\opt{A}{2 N/m}
\opt{B}{20 N/m}
\opt{C}{2 N/cm}
\opt{D}{2 N/mm}
\end{qoptions}

\sect{Section-I(ii)}

\qn{5} The figure shows the variation of photo current with anode potential for a photosensitive surface for three different radiations. Let $I_a$, $I_b$ and $I_c$ be the intensities and $f_a$, $f_b$ and $f_c$ be the frequencies for the curves $a$, $b$ and $c$ respectively. Choose correct options
\qfig{p2_phy_s2_q05_1}
\begin{qoptions}
\opt{A}{$f_a = f_b$}
\opt{B}{$I_a < I_b$}
\opt{C}{$f_c < f_b$}
\opt{D}{$I_c > I_b$}
\end{qoptions}

\qn{6} The graph shown by the side shows the variation of potential energy $\phi$ of a proton with its distance `r' from a fixed sodium nucleus, as it approaches the nucleus, placed at origin O. Then the portion:
\qfig{p2_phy_s2_q06_1}
\begin{qoptions}
\opt{A}{AB indicates nuclear repulsion}
\opt{B}{AB indicates electrostatic repulsion}
\opt{C}{BC indicates nuclear attraction}
\opt{D}{BC represents electrostatic interaction}
\end{qoptions}

\qn{7} Two bodies A and B have thermal emissivities of 0.01 and 0.81 respectively. The outer surface areas of the two bodies are the same. The two bodies emit total radiant power of the same rate. The wavelength $\lambda_B$ corresponding to maximum spectral radiancy in the radiation from B is shifted from the wavelength corresponding to maximum spectral radiancy in the radiation from A, by 1.00 $\mu$m. If the temperature of A is 5802 K:
\begin{qoptions}
\opt{A}{the temperature of B is 1934 K}
\opt{B}{$\lambda_B = 1.5\ \mu$m}
\opt{C}{the temperature of B is 11604 K}
\opt{D}{the temperature of B is 2901 K}
\end{qoptions}

\qn{8} The water equivalent of a copper calorimeter is 4.5. If the specific heat of copper is 0.09 cal/g/$^\circ$C,
\begin{qoptions}
\opt{A}{mass of the calorimeter is 0.5 kg}
\opt{B}{thermal capacity of the calorimeter is 4.5 cal/$^\circ$C}
\opt{C}{heat required to raise the temperature of the calorimeter by 8$^\circ$C will be 36 cal}
\opt{D}{heat required to melt 15 g of ice placed in the calorimeter will be 1200 cal}
\end{qoptions}

\qn{9} Two rods of different materials are placed between massive walls as shown in figure. The cross section of the rods is $A$, their moduli of elasticity are $E_1$ and $E_2$ respectively. If rods are heated by $T$ degrees, then (coefficients of linear expansion of material of rods are $\alpha_1$ and $\alpha_2$ respectively)
\qfig{p2_phy_s2_q09_1}
\begin{qoptions}
\opt{A}{the strain produced in both rods is equal.}
\opt{B}{the stress produced in both the rods is equal.}
\opt{C}{the force exerted by one rod on the other is $\dfrac{\alpha_1\ell_1+\alpha_2\ell_2}{\frac{\ell_1}{E_1}+\frac{\ell_2}{E_2}}AT$.}
\opt{D}{the force exerted by one rod on the other is $\dfrac{(E_1-E_2)^2A}{E_1+E_2}$.}
\end{qoptions}

\qn{10} A mixture of ideal gas containing 5 moles of monatomic gas and 1 mole of rigid diatomic gas is initially at pressure $P_0$, volume $V_0$ and temperature $T_0$. If the gas mixture is adiabatically compressed to a volume $V_0/4$, then the correct statement(s) is/are, (Give $2^{1.2} = 2.3$ ; $2^{3.2} = 9.2$; $R$ is gas constant)
\begin{qoptions}
\opt{A}{The final pressure of the gas mixture after compression is in between $9P_0$ and $10P_0$}
\opt{B}{The average kinetic energy of the gas mixture after compression is in between $18RT_0$ and $19RT_0$}
\opt{C}{The work $|W|$ done during the process is $13RT_0$}
\opt{D}{Adiabatic constant of the gas mixture is 1.6}
\end{qoptions}

\sect{Section-III}

\qn{1} An aluminium vessel of mass 0.5 kg contains 0.2kg of water at 20$^\circ$C. A block of iron of mass 0.2 kg at 100$^\circ$C is gently put it to water. If the equilibrium temperature of mixture is T$^\circ$C, fill the value of $\dfrac{T}{5}$. $S_{\text{iron}} = 470$ J kg$^{-1}$k$^{-1}$; $S_{\text{A}\ell} = 1140$ J kg$^{-1}$k$^{-1}$; $S_{\text{water}} = 4200$ J kg$^{-1}$k$^{-1}$

\qn{2} The cylindrical filament of an incandescent lamp of power 64 W is made of Tungsten. The operation temperature of the lamp is 2000K. Consider the filament a black body and find its radius (in mm).\\
\relax[Given : $\sigma = 6\times10^{-8}$ W/m$^2$ \& length of filament is $\dfrac{10}{3\pi}$ cm]

\qn{3} An ideal monoatomic gas undergoes a cyclic process represented by the PV graph shown in the figure. If the efficiency of this process is $\dfrac{\alpha}{\beta}$, then find the value of $\dfrac{\alpha+\beta}{9}$.
\par\nobreak
\qfig{p2_phy_s3_q03_1}

\qn{4} One mole of a diatomic gas is taken around a cyclic process shown in the figure. The process BC is adiabatic expansion and the process CA is isothermal compression. If the volume of the gas at C is given by $V_C = (2)^n V_0$, then the value of n is
\par\nobreak
\qfig{p2_phy_s3_q04_1}

\qn{5} An electron in a hydrogen atom in excited state n reaches the ground state by emitting three photons of frequency in ratio 5 : 72 : 243. Wavelength corresponding to one of the transitions belongs to visible region. Find the value of `n'.

\qn{6} Radiation from a hydrogen discharge tube (energy of photons $\le$ 13.6 eV) goes through a filter which transmits only waves of wavelength greater than 4400 \AA\ and is incident on a metal of work function 2.0 eV. If stopping potential is $11n\times10^{-2}$ volts. Find the value of `n'. (take $hc$ = 12400 eV\AA)

\qn{7} When the voltage applied to an X-ray tube increased from $V_1 = 15.5$ KV to $V_2 = 31$ KV, the wavelength interval between the K$_\alpha$ line and the short wavelength cut-off of the continuous X-ray spectrum increases by a factor of 1.3. Find the atomic number of the element of the target. If your answer is N fill value of N/13. (Take : $hc$ = 1240 eV nm and $R = 1\times10^7$ /m)

\qn{8} A parallel beam of intensity 3000 W/m$^2$, $\lambda$ = 100 nm having cross sectional area 100 cm$^2$ is incident normally on a glass slab, placed in air. The slab reflects 40 \% of light, transmits 40 \% of light and absorbs 20 \% of light incident on it. The force experienced by slab due to incident radiation beam will be $n\times10^{-7}$ N. Find the value of n.

\partsec{2}{CHEMISTRY}
\sect{Section-I(i)}

\qn{1} Calculate the Van't hoff factor when 0.1 mole of NH$_4$Cl is dissolved in 1 liter water solution. The degree of dissociation of NH$_4$Cl is 0.8 and degree of hydrolysis of NH$_4^+$ ion is 0.1 :
\begin{qoptions}
\opt{A}{0.8}
\opt{B}{1.8}
\opt{C}{1.88}
\opt{D}{2.68}
\end{qoptions}

\qn{2} In XeO$_3$F$_2$ number of $\pi$-bonds and total number of lone pairs respectively are :-
\begin{qoptions}
\opt{A}{3, 12}
\opt{B}{3, 6}
\opt{C}{2, 12}
\opt{D}{2, 10}
\end{qoptions}

\qn{3} Which intermolecular force is most responsible in allowing Kr to liquify?
\begin{qoptions}
\opt{A}{Ionic bond}
\opt{B}{Instantaneous dipole-induced dipole force}
\opt{C}{Dipole-dipole force}
\opt{D}{Ion-dipole force}
\end{qoptions}

\qn{4}\par\nobreak
\qfig{p2_chem_s1_q04_1}
In the above conversion of compound (X) to product (Y), the sequence of reagents to be used will be:
\begin{qoptions}
\opt{A}{(i) Br$_2$, Fe (ii) Fe, H$^+$ (iii) LiAlH$_4$}
\opt{B}{(i) Br$_2$(aq) (ii) LiAlH$_4$ (iii) H$_3$O$^+$}
\opt{C}{(i) Fe, H$^+$ (ii) Br$_2$(aq) (iii) HNO$_2$ (iv) CuBr}
\opt{D}{(i) Fe, H$^+$ (ii) Br$_2$(aq) (iii) HNO$_2$ (iv) H$_3$PO$_2$}
\end{qoptions}

\sect{Section-I(ii)}

\qn{5} Select the correct statement(s), related with the following aqueous solutions :\\
(Assume $\alpha = 1$ if required)
\begin{qoptions}
\opt{A}{0.1 m CaCl$_2$ has higher boiling point than 0.1 m NaCl solution.}
\opt{B}{0.05 m Al$_2$(SO$_4$)$_3$ has higher freezing point than 0.1 m K$_3$[Fe(CN)$_6$] solution.}
\opt{C}{0.1 M glucose (C$_6$H$_{12}$O$_6$) exert higher osmotic pressure than 0.08 m CH$_3$COOH (25\% dissociated).}
\opt{D}{Vapour pressure of 0.05 M urea solution is greater than that of 0.05 M KCl solution}
\end{qoptions}

\qn{6} A colloid of ferric oxide is prepared by mixing FeCl$_3$ in excess of NaOH solution. Select the correct statement(s) for this colloid :
\begin{qoptions}
\opt{A}{In electro-osmosis, the sol is coagulated at cathode.}
\opt{B}{It can be coagulated by anions.}
\opt{C}{It is mutual coagulated by addition of colloid of As$_2$S$_3$.}
\opt{D}{It is negatively charged sol.}
\end{qoptions}

\qn{7} MgSO$_4$.7H$_2$O (Epsom salt) have following type of force of attraction.
\begin{qoptions}
\opt{A}{Ionic bond}
\opt{B}{Covalent bond}
\opt{C}{Coordinate bond}
\opt{D}{Hydrogen bond}
\end{qoptions}

\qn{8} Which of the following order is/are correct?
\begin{qoptions}
\opt{A}{Bond energy : Cl$_2$ $>$ Br$_2$ $>$ F$_2$ $>$ I$_2$}
\opt{B}{Bond strength : $-\underset{|}{\mathrm N}-\underset{|}{\mathrm N}- \;>\; -\underset{|}{\mathrm P}-\underset{|}{\mathrm P}-$}
\opt{C}{Number of $\pi$-bonds : HOCl $<$ HClO$_2$ $<$ HClO$_3$ $<$ HClO$_4$}
\opt{D}{Bond angle : NO$_2^+$ $>$ NO$_2$ $>$ NO$_2^-$}
\end{qoptions}

\qn{9} Which of the following method(s) is/are not used to prepare p-bromo aniline as major product :
\begin{qoptions}
\opt{A}{\qfigi{p2_chem_s2_q09_1}}
\opt{B}{\qfigi{p2_chem_s2_q09_2}}
\opt{C}{\qfigi{p2_chem_s2_q09_4}}
\opt{D}{\qfigi{p2_chem_s2_q09_3}}
\end{qoptions}

\qn{10} Which of the following statements are correct about the structure of starch?
\begin{qoptions}
\opt{A}{It is a polymer of $\alpha$-D (+) glucose linked through $\alpha$-glycosidic bond.}
\opt{B}{It consists of water soluble long unbranched chain of amylose.}
\opt{C}{It consists of water insoluble branched amylopectin having C$_1$--C$_4$ and C$_1$--C$_6$ glycosidic linkage between $\alpha$-D glucose.}
\opt{D}{In animal body it is stored in the form of glycogen, also known as animal starch.}
\end{qoptions}

\sect{Section-III}

\qn{1} At 17$^\circ$C, the osmotic pressure of 1 litre sucrose solution is 580 torr. The solution is diluted and the temperature is raised to 57$^\circ$C, when the osmotic pressure is found to be 165 torr. The volume of diluted solution is \underline{\hspace{2em}} litre.

\qn{2} In an adsorption experiment,
\par\nobreak
\qfig{p2_chem_s3_q02_1}
Calculate the mass of gas (in gram) adsorbed over 2 gram charcoal under a pressure of 2 atm.

\qn{3} Henry's law constant for CO$_2$ in water is $2.0\times10^8$ pascal at 27$^\circ$C. The volume of CO$_2$ gas in 900 g soda water when packed under $8.0\times10^5$ pascal CO$_2$ pressure at same temperature is \underline{\hspace{2em}} $\times\,10^2$ millilitre.\\
(1 atm = $10^5$ pascal, $R$ = 0.08 atm-L/mole-K)

\qn{4} On the basis of molecular orbital theory, sum of $\pi$-bonds present in B$_2$, C$_2$ and O$_2$ is :-

\qn{5} $a + b + c$ = ?\\
where\\
$a$ = Number of 90 degrees F-I-F angles in IF$_5$\\
$b$ = Basicity of H$_3$PO$_2$\\
$c$ = Number of banana bonds in B$_2$H$_6$

\qn{6} The total number of distinct naturally occurring amino acids obtained by complete acidic hydrolysis of the peptide shown below is
\par\nobreak
\qfig{p2_chem_s3_q06_1}

\qn{7} Among the following number of condensation polymers are ?\\
Nylon-6, Buna-N, Buna-S, Nylon-6,6, Nylon-6,10, PVC, Polystyrene, Teflon.

\qn{8} In an estimation of bromine by Carius method, 1.6 g of an organic compound gave 1.88 g of AgBr. The mass percentage of bromine in the compound is \underline{\hspace{3em}}.\\
(Atomic mass, Ag=108, Br = 80 g mol$^{-1}$)\\
\textbf{\textit{Fill your answer as sum of digits (excluding decimal places) till you get the single digit answer.}}

\partsec{3}{MATHEMATICS}
\sect{Section-I(i)}

\qn{1} A \& B are two points on $\sqrt3 x+y+1=0$ such that $AB$ is 3 and area of $\Delta ABC$ is 9. If locus of $C$ in $xy$-plane is a pair of two parallel lines, then distance between them is -
\begin{qoptions}
\opt{A}{1}
\opt{B}{3}
\opt{C}{6}
\opt{D}{12}
\end{qoptions}

\qn{2} The $\omega$ is a non - real cube root of unity then value of $\dfrac{6\left(5+\omega^2\right)}{(5+\omega)^{-1}}+\dfrac{7\left(6+\omega^2\right)}{(6+\omega)^{-1}}+\ldots\ldots+\dfrac{21\left(20+\omega^2\right)}{(20+\omega)^{-1}}$ is:
\begin{qoptions}
\opt{A}{42015}
\opt{B}{44016}
\opt{C}{43015}
\opt{D}{45012}
\end{qoptions}

\qn{3} The locus of the centres of the circles which cut the circles $x^2+y^2+4x-6y+9=0$ and $x^2+y^2-5x+4y-2=0$ orthogonally is--
\begin{qoptions}
\opt{A}{$9x+10y-7=0$}
\opt{B}{$x-y+2=0$}
\opt{C}{$9x-10y+11=0$}
\opt{D}{None of these}
\end{qoptions}

\qn{4} If $a$ and $b\in R$ and satisfy the relation $4a^2+3b^2-4ab-4b+2\le 0$ then value of determinant $\begin{vmatrix}4&4&4\\a&b&2\\a^2&b^2&4\end{vmatrix}$ is equal to -
\begin{qoptions}
\opt{A}{2}
\opt{B}{3}
\opt{C}{$-3$}
\opt{D}{0}
\end{qoptions}

\sect{Section-I(ii)}

\qn{5} Let $B=[b_{ij}]_{n\times n}$ be a square matrix of order $n$ such that $b_{ij}=\begin{cases}0 & ,\ i\ne j\\[2pt] \dfrac{4+i+j}{2} & ,\ i=j\end{cases}$
\begin{qoptions}
\opt{A}{$\displaystyle\lim_{m\to\infty}\left(\dfrac{\sum\limits_{n=1}^{m}\mathrm{trace}(B)}{m^3}\right)=\dfrac16$}
\opt{B}{$\displaystyle\sum_{n=1}^{10}\mathrm{trace}(B)=330$}
\opt{C}{$\displaystyle\sum_{n=1}^{10}\mathrm{trace}(B)=320$}
\opt{D}{$\displaystyle\lim_{m\to\infty}\left(\dfrac{\sum\limits_{n=1}^{m}\mathrm{trace}(B)}{m^3}\right)=\dfrac13$}
\end{qoptions}

\qn{6} Let $a,b,c$ be distinct complex numbers with $|a|=|b|=|c|=1$ and $z_1,z_2$ be the roots of the equation $az^2+bz+c=0$ with $|z_1|=1$. Also $P$ and $Q$ are the points representing the complex numbers $z_1$ and $z_2$ respectively in the complex plane with $\angle POQ=\theta$ (where $O$ being the origin) then which of the following is/are \textbf{correct}?
\begin{qoptions}
\opt{A}{$b^2=ac$}
\opt{B}{$\theta=\dfrac{2\pi}{3}$}
\opt{C}{$PQ=\sqrt3$}
\opt{D}{$|z_1+z_2|=1$}
\end{qoptions}

\qn{7} Consider lines $L_1: x+2y-3=0$ and $L_2: x+y-5=0$ intersecting each other at point $A$. From point $B(3,0)$, a line is drawn intersecting $L_2$ at $D$ such that $BD$ is minimum and hence a parallelogram $ABCD$ is completed. Then -
\begin{qoptions}
\opt{A}{Area of parallelogram $ABCD$ is 6 sq. units}
\opt{B}{Area of triangle $BCD$ is 4 sq. units}
\opt{C}{Area of triangle $BCD$ is 3 sq. units}
\opt{D}{Coordinates of vertex $C$ is $(0,3)$}
\end{qoptions}

\qn{8} Given $p,q\in\{0,1,2,3,\ldots\ldots..,19,20\}$, consider the system of equations\\
$x+y+z=4$\\
$2x+y+3z=6$\\
$x+2y+pz=q$\\
Let\\
$L$ denotes number of ordered pairs $(p,q)$ so that system has a unique solution\\
$M$ denotes number of ordered pairs $(p,q)$ so that system has no solution\\
$N$ denotes number of ordered pairs $(p,q)$ so that system has infinite solutions,\\
then :
\begin{qoptions}
\opt{A}{$L=420$}
\opt{B}{$M=20$}
\opt{C}{$N=2$}
\opt{D}{$\dfrac{L+M+N}{21}=21$}
\end{qoptions}

\qn{9} If direct common tangents of circle $S_1\equiv x^2+y^2-2x-4y-4=0$ and $S_2\equiv x^2+y^2-8x-12y+36=0$ touches circle $S_1=0$ at $A$ and $B$ and $S_3=0$ is equation of circle whose diametrically opposite end points are $A$ and $B$, then
\begin{qoptions}
\opt{A}{equation of $S_3$ is $25x^2+25y^2-32x-76y-148=0$}
\opt{B}{equation of $S_3$ is $x^2+y^2+7x+8y-28=0$}
\opt{C}{length of common chord of $S_1=0$ and $S_2=0$ is $\dfrac{24}{5}$}
\opt{D}{length of common chord of $S_1=0$ and $S_2=0$ is $\dfrac{18}{5}$}
\end{qoptions}

\qn{10} Let $\Delta=\begin{vmatrix}a&b&c\\a-b&b-c&c-a\\b+c&c+a&a+b\end{vmatrix}$.\\
Which option is/are always correct
\begin{qoptions}
\opt{A}{If $a,b,c$ are unequal positive real numbers, then $\Delta$ is always positive}
\opt{B}{If $\Delta=0$, then $a^3+b^3+c^3=3abc$}
\opt{C}{If $a,b,c$ are unequal negative real numbers, then $\Delta$ is always negative}
\opt{D}{If $\Delta=0$, then $a^2+b^2+c^2=ab+bc+ca$}
\end{qoptions}

\sect{Section-III}

\qn{1} If $\sqrt{x+3-4\sqrt{x-1}}+\sqrt{(x+8)-6\sqrt{x-1}}=1$, then the number of integers satisfying above equation will be equal to

\qn{2} There are two possible values of $A$ (say $A_1$ \& $A_2$) in the solution of matrix equation $\begin{bmatrix}2A+1&-5\\-4&A\end{bmatrix}^{-1}\begin{bmatrix}A-5&B\\2A-2&C\end{bmatrix}=\begin{bmatrix}14&D\\E&F\end{bmatrix}$ then find $|-27(A_1+A_2)|$ is

\qn{3} Let $z$ and $w$ be two complex numbers such that $w=z\bar z-2z+2$, $\left|\dfrac{z+i}{z-3i}\right|=1$ and $\mathrm{Re}(w)$ has minimum value. Then, the minimum value of $n\in N$ for which $w^n$ is real, is equal to \underline{\hspace{2.5em}}.

\qn{4} The sum of 162$^{\text{th}}$ power of the roots of the equation $x^3-2x^2+2x-1=0$ is

\qn{5} A line which makes an acute angle `$\theta$' with positive direction of x-axis is drawn from $P(5,4)$ to meet the line $x=-7$ at $Q$ and $y=9$ at $R$ then $\dfrac{144}{PQ^2}+\dfrac{25}{PR^2}$ is

\qn{6} If $A=\begin{bmatrix}0&2\\0&0\end{bmatrix}$ \& $(I+A)^{20}-19A=\begin{bmatrix}\alpha&\beta\\\gamma&\delta\end{bmatrix}$, then value of $\alpha+\beta+\gamma+\delta$ is equal to\\
(Here $I$ is $2\times2$ unit matrix)

\qn{7} If $a=\max\{(x+2)^2+(y-3)^2\}$ and $b=\min\{(x+2)^2+(y-3)^2\}$ where $x,y$ satisfying $x^2+y^2+8x-10y-40=0$, then value of $\dfrac{a-b}{8\sqrt2}$ is

\qn{8} If equation $3x+2y-z=4$, $x+y-z=\lambda$ and $2x-\mu y+2z=3$ have infinite number of simultaneous solutions then the value of $4(\lambda+\mu)$ is
\end{multicols}
\end{document}
